% !TeX root = ../sections/02_exercises.tex
\begin{exercise}
	数列 $\{a_n\}_{n\in\mathbb{N}}$ と $\{b_n\}_{n\in\mathbb{N}}$ が収束して
	\[
	\lim_{n\to\infty} a_n=\alpha,
	\qquad
	\lim_{n\to\infty} b_n=\beta
	\]
	であり，$\alpha<\beta$ が成立するとき，
	ある番号 $N\in\mathbb{N}$ が存在して，
	$n\geq N$ ならば
	\[
	a_n<b_n
	\]
	であることを証明せよ。
\end{exercise}
\begin{background}
	この問題では，極限値の間にある大小関係
	\[
	\alpha<\beta
	\]
	を，十分大きな番号以降の数列の大小関係
	\[
	a_n<b_n
	\]
	へ移すことを考える。
	
	\begin{definition}[数列の収束]
		数列 $\{a_n\}$ が実数 $\alpha$ に収束するとは，
		任意の $\epsilon>0$ に対して，ある自然数 $N$ が存在して，
		すべての $n\geq N$ に対して
		\[
		|a_n-\alpha|<\epsilon
		\]
		が成り立つことをいう。
	\end{definition}
	
	\begin{remark}
		$a_n\to\alpha$ であるから，十分大きな $n$ に対して
		\[
		a_n<\alpha+\epsilon
		\]
		が成り立つ。
		
		また，$b_n\to\beta$ であるから，十分大きな $n$ に対して
		\[
		\beta-\epsilon<b_n
		\]
		が成り立つ。
	\end{remark}
	
	\subsection{極限値の間隔を数列へ反映する準備}
	
	\begin{lemma}[有限個の条件を同時に満たす番号の取り方]
		自然数 $N_1,N_2$ に対して
		\[
		N=\max\{N_1,N_2\}
		\]
		とおけば，$n\geq N$ のとき
		\[
		n\geq N_1,
		\qquad
		n\geq N_2
		\]
		が同時に成り立つ。
	\end{lemma}
	
	\begin{remark}
		$\alpha<\beta$ であるから，
		\[
		\beta-\alpha>0
		\]
		である。この正の差を利用して，
		$a_n$ を $\alpha$ の近くに，$b_n$ を $\beta$ の近くに
		同時に近づければ，
		最終的に
		\[
		a_n<b_n
		\]
		を得ることができる。
	\end{remark}
\end{background}
\begin{strategy}
	示したいことは，ある自然数 $N$ が存在して，
	すべての $n\geq N$ に対して
	\[
	a_n<b_n
	\]
	となることである。
	
	仮定より
	\[
	\alpha<\beta
	\]
	である。したがって
	\[
	\beta-\alpha>0
	\]
	である。
	
	この正の差を利用するために，
	\[
	\epsilon=\frac{\beta-\alpha}{3}
	\]
	とおく。
	
	$a_n\to\alpha$ であるから，この $\epsilon>0$ に対して，
	ある自然数 $N_1$ が存在して，すべての $n\geq N_1$ に対して
	\[
	|a_n-\alpha|<\epsilon
	\]
	が成り立つ。
	
	この不等式から
	\[
	-\epsilon<a_n-\alpha<\epsilon
	\]
	である。特に
	\[
	a_n<\alpha+\epsilon
	\]
	を得る。
	
	また，$b_n\to\beta$ であるから，この $\epsilon>0$ に対して，
	ある自然数 $N_2$ が存在して，すべての $n\geq N_2$ に対して
	\[
	|b_n-\beta|<\epsilon
	\]
	が成り立つ。
	
	この不等式から
	\[
	-\epsilon<b_n-\beta<\epsilon
	\]
	である。特に
	\[
	\beta-\epsilon<b_n
	\]
	を得る。
	
	ここで，$\epsilon=(\beta-\alpha)/3$ と選んでいるので，
	\[
	\alpha+\epsilon<\beta-\epsilon
	\]
	が成り立つ。実際，
	\[
	(\beta-\epsilon)-(\alpha+\epsilon)
	=
	\beta-\alpha-2\epsilon
	\]
	であり，
	\[
	\beta-\alpha-2\epsilon
	=
	\beta-\alpha-\frac{2}{3}(\beta-\alpha)
	=
	\frac{1}{3}(\beta-\alpha)>0
	\]
	である。
	
	そこで
	\[
	N=\max\{N_1,N_2\}
	\]
	とおく。
	
	このとき，$n\geq N$ ならば
	\[
	n\geq N_1,
	\qquad
	n\geq N_2
	\]
	であるから，
	\[
	a_n<\alpha+\epsilon
	\]
	かつ
	\[
	\beta-\epsilon<b_n
	\]
	が同時に成り立つ。
	
	したがって
	\[
	a_n<\alpha+\epsilon<\beta-\epsilon<b_n
	\]
	となる。
	
	よって
	\[
	a_n<b_n
	\]
	が成り立つ。
	
	以上より，ある自然数 $N$ が存在して，すべての $n\geq N$ に対して
	\[
	a_n<b_n
	\]
	となる。
\end{strategy}